\subsubsection*{Численный пример (ДФП)}

\begin{idea}[title={Откуда этот материал}]
Прогон ДФП на числах; на лекции формулы были выписаны без примера. Обозначения
стандартные: $s_k = x_{k+1} - x_k$, $y_k = \nabla f(x_{k+1}) - \nabla f(x_k)$,
нумерация точек с нуля.
\end{idea}

\begin{example}\label{ex:dfp-run}
\[
  f(x_1, x_2) = x_1^2 + 2x_2^2 \ \longrightarrow\ \min ,
  \qquad
  \nabla f(x) = \begin{pmatrix} 2x_1 \\ 4x_2 \end{pmatrix} ,
  \qquad
  x_0 = \begin{pmatrix} 2 \\ 1 \end{pmatrix} , \quad A_0 = E .
\]
Шаг выбирается точной одномерной минимизацией.

\medskip
\noindent\textbf{Итерация 0.}
\[
  g_0 = \nabla f(x_0) = \begin{pmatrix} 4 \\ 4 \end{pmatrix} ,
  \qquad
  p_0 = -A_0 g_0 = \begin{pmatrix} -4 \\ -4 \end{pmatrix}
\]
(первый шаг — обычный антиградиентный, потому что $A_0 = E$). Одномерная
задача:
\[
  \varphi(\alpha) = f(x_0 + \alpha p_0)
    = (2 - 4\alpha)^2 + 2(1 - 4\alpha)^2
    = 6 - 32\alpha + 48\alpha^2 ,
\]
\[
  \varphi'(\alpha) = -32 + 96\alpha = 0
  \quad\Longrightarrow\quad
  \alpha_0 = \tfrac13 ,
\]
\[
  x_1 = x_0 + \alpha_0 p_0 = \begin{pmatrix} 2/3 \\ -1/3 \end{pmatrix} ,
  \qquad
  f(x_1) = \tfrac49 + \tfrac29 = \tfrac23
  \qquad (6 \to \tfrac23) .
\]

\medskip
\noindent\textbf{Пересчёт матрицы.}
\[
  g_1 = \nabla f(x_1) = \begin{pmatrix} 4/3 \\ -4/3 \end{pmatrix} ,
  \qquad
  s_0 = x_1 - x_0 = \begin{pmatrix} -4/3 \\ -4/3 \end{pmatrix} ,
  \qquad
  y_0 = g_1 - g_0 = \begin{pmatrix} -8/3 \\ -16/3 \end{pmatrix} .
\]
Условие кривизны:
\[
  s_0\T y_0 = \tfrac{32}{9} + \tfrac{64}{9} = \tfrac{96}{9} = \tfrac{32}{3} > 0 .
\]
Поскольку $A_0 = E$, имеем $A_0 y_0 = y_0$ и
$y_0\T A_0 y_0 = \tfrac{64}{9} + \tfrac{256}{9} = \tfrac{320}{9}$. Слагаемые
формулы ДФП:
\[
  \frac{s_0 s_0\T}{s_0\T y_0}
    = \frac{16/9}{32/3}\begin{pmatrix} 1 & 1 \\ 1 & 1 \end{pmatrix}
    = \frac16 \begin{pmatrix} 1 & 1 \\ 1 & 1 \end{pmatrix} ,
\]
\[
  \frac{A_0 y_0 y_0\T A_0}{y_0\T A_0 y_0}
    = \frac{9}{320}\begin{pmatrix} 64/9 & 128/9 \\ 128/9 & 256/9 \end{pmatrix}
    = \begin{pmatrix} 1/5 & 2/5 \\ 2/5 & 4/5 \end{pmatrix} .
\]
Итого
\[
  A_1 = A_0 + \frac{s_0 s_0\T}{s_0\T y_0}
            - \frac{A_0 y_0 y_0\T A_0}{y_0\T A_0 y_0}
      = \begin{pmatrix} 29/30 & -7/30 \\[2pt] -7/30 & 11/30 \end{pmatrix} .
\]
Проверка секущего условия (теорема~\ref{thm:dfp-secant}):
\[
  A_1 y_0
   = \begin{pmatrix} 29/30 & -7/30 \\ -7/30 & 11/30 \end{pmatrix}
     \begin{pmatrix} -8/3 \\ -16/3 \end{pmatrix}
   = \begin{pmatrix} -4/3 \\ -4/3 \end{pmatrix} = s_0 .
\]

\medskip
\noindent\textbf{Итерация 1.}
\[
  A_1 g_1
   = \begin{pmatrix} 29/30 & -7/30 \\ -7/30 & 11/30 \end{pmatrix}
     \begin{pmatrix} 4/3 \\ -4/3 \end{pmatrix}
   = \begin{pmatrix} 8/5 \\ -4/5 \end{pmatrix} ,
  \qquad
  p_1 = -A_1 g_1 = \begin{pmatrix} -8/5 \\ 4/5 \end{pmatrix} .
\]
\[
  \varphi(\alpha) = f(x_1 + \alpha p_1)
    = \Big(\tfrac23 - \tfrac85\alpha\Big)^2
      + 2\Big(-\tfrac13 + \tfrac45\alpha\Big)^2
    = \tfrac23 - \tfrac{16}{5}\alpha + \tfrac{96}{25}\alpha^2 ,
\]
\[
  \varphi'(\alpha) = -\tfrac{16}{5} + \tfrac{192}{25}\alpha = 0
  \quad\Longrightarrow\quad
  \alpha_1 = \tfrac{5}{12} ,
\]
\[
  x_2 = x_1 + \alpha_1 p_1
      = \begin{pmatrix} 2/3 - 2/3 \\ -1/3 + 1/3 \end{pmatrix}
      = \begin{pmatrix} 0 \\ 0 \end{pmatrix} = x_* ,
  \qquad f(x_2) = 0 .
\]
Значения: $6 \to \tfrac23 \to 0$.
\end{example}

\begin{idea}
Точное попадание в минимум за две итерации — не случайность. Для
квадратичной функции $n$ переменных квазиньютоновский метод с \emph{точным}
одномерным поиском завершается не более чем за $n$ шагов: накопленные пары
$(s_k, y_k)$ полностью восстанавливают гессиан, и после $n$-й итерации
$A_n = H^{-1}$. Здесь $n = 2$, поэтому хватило двух шагов. Проверить это можно
прямо: $H = \begin{pmatrix} 2 & 0 \\ 0 & 4 \end{pmatrix}$,
$H^{-1} = \begin{pmatrix} 1/2 & 0 \\ 0 & 1/4 \end{pmatrix}$ — и именно к этой
матрице сошлась бы $A_2$, если бы её потребовалось строить.

Сравните с примером~\ref{ex:one-dim-run} по духу: там тоже использовались
только значения функции вдоль прямой, но здесь направление каждый раз
корректируется накопленной кривизной, а не берётся по антиградиенту.
\end{idea}
